\documentclass[11pt,a4paper]{article}
\usepackage[utf8]{inputenc}
\usepackage[T1]{fontenc}
\usepackage[portuguese]{babel}
\usepackage[a4paper,margin=2.2cm]{geometry}
\usepackage{array}
\usepackage{booktabs}
\usepackage{longtable}
\usepackage{tabularx}
\usepackage{xcolor}
\usepackage{enumitem}
\usepackage{hyperref}
\usepackage{titlesec}
\usepackage{fancyhdr}

\hypersetup{
  colorlinks=true,
  linkcolor=blue!55!black,
  urlcolor=blue!55!black,
  pdftitle={Regras, Permissões e Vistas por Perfil},
  pdfauthor={SMARTER HUB}
}

\pagestyle{fancy}
\fancyhf{}
\lhead{SMARTER HUB}
\rhead{Regras por perfil}
\cfoot{\thepage}

\definecolor{titleblue}{HTML}{123F86}
\definecolor{softblue}{HTML}{EAF2FF}
\definecolor{softgray}{HTML}{F7F9FC}
\definecolor{accent}{HTML}{2F7DF1}

\titleformat{\section}{\large\bfseries\color{titleblue}}{}{0pt}{}
\titleformat{\subsection}{\normalsize\bfseries\color{titleblue}}{}{0pt}{}

\renewcommand{\arraystretch}{1.25}
\setlist[itemize]{leftmargin=1.2em, itemsep=0.25em, topsep=0.25em}
\setlist[enumerate]{leftmargin=1.5em, itemsep=0.25em, topsep=0.25em}

\newcolumntype{Y}{>{\raggedright\arraybackslash}X}

\begin{document}

\begin{center}
  {\LARGE\bfseries\color{titleblue} Regras, Permissões e Vistas por Perfil}

  \vspace{0.4em}
  {\normalsize Documento de referência funcional do portal}
\end{center}

\vspace{0.5em}
\noindent\colorbox{softblue}{\parbox{\dimexpr\linewidth-2\fboxsep\relax}{
\textbf{Objetivo.} Este documento resume o comportamento esperado do sistema para menus, visibilidade, aprovações e edição. Sempre que houver dúvida funcional, esta é a base a seguir.
}}

\section{Regras base}
\begin{itemize}
  \item Um utilizador pode pertencer a mais do que uma equipa.
  \item Um utilizador pode ser chefe de uma equipa e participante noutra.
  \item O pedido de férias ou ausências deve indicar a equipa ou subequipa de contexto.
  \item A aprovação segue a cadeia dessa equipa ou subequipa.
  \item Quando o pedido nasce numa subequipa, a cadeia deve herdar os níveis da equipa mãe até fechar a hierarquia.
  \item Em multi-chefia, todos os aprovadores relevantes da linha devem aprovar.
  \item Em Portugal, o admin pode aprovar por exceção e substituir a cadeia normal.
  \item Um pedido aprovado pode ser alterado por versionamento, não por edição direta do registo final.
  \item Deve existir histórico de decisões, versões e motivos.
  \item Deve evitar-se ter cerca de um terço da equipa em férias em simultâneo.
  \item Pedidos que atravessam o fim/início de ano devem aparecer nas vistas do(s) ano(s) respetivo(s).
  \item Ausências podem usar horas e podem ter anexos pendentes sem prazo fixo.
  \item Férias podem ser pedidas a meio dia.
\end{itemize}

\section{Matriz por perfil}

\rowcolors{2}{softgray}{white}
\begin{longtable}{p{2.4cm} p{4.4cm} p{3.2cm} p{3.0cm} p{3.2cm}}
\toprule
\textbf{Perfil} & \textbf{Pode ver} & \textbf{Pode criar} & \textbf{Pode aprovar} & \textbf{Pode editar} \\
\midrule
\endfirsthead
\toprule
\textbf{Perfil} & \textbf{Pode ver} & \textbf{Pode criar} & \textbf{Pode aprovar} & \textbf{Pode editar} \\
\midrule
\endhead

Colaborador & A sua ficha, notificações, formações, férias e ausências, a sua equipa e as férias dos membros dessa equipa & Pedido de férias, ausências, pedido de alteração da ficha, submissão de comprovativos, autoatribuição de formação quando permitido & Não aprova pedidos operacionais & A própria ficha, pedidos pendentes por versionamento, comprovativos, pedido de alteração de carga horária \\

Manager & As equipas sob a sua responsabilidade, membros dessas equipas, férias e ausências dos subordinados, formações atribuídas, aprovações pendentes do seu nível & Atribuir formações, gerir membros dentro da sua esfera autorizada, criar pedidos no âmbito da sua gestão & Aprova férias e ausências da sua linha de chefia, conforme a cadeia da equipa e contexto & Pode editar o que estiver na sua esfera de gestão, não dados sensíveis globais nem equipas fora da sua responsabilidade \\

Coordenador & Equipas sob a sua responsabilidade, férias e ausências das chefias e direção sob supervisão, aprovações pendentes do seu nível & Atribuir formações quando autorizado, apoiar gestão de equipas dentro do seu âmbito & Aprova os níveis superiores da hierarquia, incluindo chefias e direção, e atua como segunda linha quando aplicável & Pode editar apenas o que estiver dentro da sua responsabilidade funcional \\

Admin & Tudo: todas as equipas, membros, férias, ausências, aprovações, fichas, formações, notificações e histórico & Criar, editar e remover equipas; gerir memberships; realocar membros; editar dados sensíveis; aprovar por exceção em PT; gerir tudo & Pode aprovar tudo, incluindo exceções de PT e substituições de cadeia quando necessário & Pode editar equipas, memberships, dados sensíveis, permissões, perfis e estrutura global \\

HR (futuro) & A definir, atualmente absorvido pelo admin & A definir & A definir & A definir \\
\bottomrule
\end{longtable}

\section{Menus por perfil}

\subsection{Colaborador}
\begin{itemize}
  \item Home
  \item A Minha Ficha
  \item Equipas
  \item Formações
  \item Férias
  \item Recibos
\end{itemize}

\subsection{Manager}
\begin{itemize}
  \item Home
  \item A Minha Ficha
  \item Equipas
  \item Aprovações
  \item Formações
  \item Férias
\end{itemize}

\subsection{Coordenador}
\begin{itemize}
  \item Home
  \item Equipas
  \item Aprovações
  \item A Minha Ficha
  \item Formações
  \item Férias
\end{itemize}

\subsection{Admin}
\begin{itemize}
  \item Home
  \item A Minha Ficha
  \item Equipas
  \item Administração
  \item Aprovações
  \item Formações
  \item Férias
  \item Recibos
  \item Notificações
\end{itemize}

\subsection{Convidado}
\begin{itemize}
  \item Home
  \item Onboarding
  \item Notificações
\end{itemize}

\section{Vistas obrigatórias}

\subsection{Colaborador}
\begin{itemize}
  \item Deve ver a sua equipa.
  \item Deve ver as férias e ausências dos membros da sua equipa.
  \item Deve conseguir pedir e editar pedidos pendentes por versionamento.
\end{itemize}

\subsection{Manager}
\begin{itemize}
  \item Deve ver as equipas sob a sua responsabilidade.
  \item Deve ver as férias e ausências dos subordinados.
  \item Deve conseguir aprovar pedidos da sua linha de chefia.
\end{itemize}

\subsection{Coordenador}
\begin{itemize}
  \item Deve ver apenas as equipas sob a sua responsabilidade.
  \item Deve aprovar a segunda linha ou os níveis superiores definidos.
\end{itemize}

\subsection{Admin}
\begin{itemize}
  \item Deve ver tudo.
  \item Deve conseguir gerir equipas, memberships, utilizadores, dados sensíveis e exceções.
\end{itemize}

\section{Estados que devem ficar claros na interface}
\begin{itemize}
  \item Pedido pendente
  \item Pedido aprovado
  \item Pedido recusado
  \item Pedido cancelado
  \item Pedido substituído por nova versão
  \item Linha de aprovação atual
  \item Equipa principal
  \item Subequipa
  \item Membro participante
  \item Membro chefe ou aprovador
  \item Pedido de meio-dia (manhã ou tarde)
\end{itemize}

\section{Observações atuais}
\begin{itemize}
  \item HR não existe como role final neste momento.
  \item O fluxo de administração de equipas e memberships deve ser feito por admin.
  \item A página de equipas deve ser comum, mas mostrar conteúdo diferente por perfil.
  \item O histórico de férias e decisões deve manter a versão anterior quando houver edição de um pedido já aprovado.
\end{itemize}

\vspace{0.5em}
\noindent\colorbox{softblue}{\parbox{\dimexpr\linewidth-2\fboxsep\relax}{
\textbf{Nota prática.} Este ficheiro é o resumo funcional fechado para orientar menus, permissões e páginas. Se houver alteração ao processo, este documento deve ser atualizado antes de mexer no código.
}}

\end{document}
